%\chapter*{Resumen}

%% Si se quiere enviar este resumen en una p\'agina a alguien, descomentar esto.

%\begin{center}
%\textbf{\large Universidad Argentina De la Empresa}

%\textbf{Facultad de Ingenier\'{\i}a y Ciencias Exactas}

%\textbf{Departamento de Tecnolog\'{\i}a Inform\'atica}

%\end{center}

%\begin{center}
%\large Ingenier\'{\i}a en Inform\'atica
%\end{center}

%\begin{center}
%\textbf{\large Título de la PFI}
%\end{center}

%\begin{center}
%Apellido, Nombres
%Apellido, Nombres
%\end{center}

%\noindent
%El envejecimiento de la población genera una demanda creciente de soluciones que permitan a los familiares mantenerse informados sobre el estado y la rutina de los adultos mayores a su cargo, especialmente cuando la distancia o las obligaciones cotidianas impiden una presencia continua. Los sistemas existentes requieren en su mayoría que el adulto mayor porte algún dispositivo o interactúe activamente con la tecnología, lo que limita su efectividad en situaciones de emergencia.

%Este proyecto propone el desarrollo de un sistema de monitoreo completamente no invasivo para adultos mayores en entornos domésticos, basado en visión artificial y aprendizaje adaptativo del comportamiento. El sistema detecta caídas, inmovilidad prolongada y actividades de la vida cotidiana, a partir del análisis de video en tiempo real mediante modelos de estimación de poses. Adicionalmente, construye de forma continua una línea base de rutina personalizada por individuo, permitiendo identificar cambios en los patrones habituales de comportamiento.

%El procesamiento se realiza localmente en una controladora, garantizando la privacidad del usuario, y solo los eventos detectados se transmiten a una aplicación móvil para la notificación a los familiares. El sistema se valida en un entorno doméstico controlado, demostrando la viabilidad del enfoque propuesto.
